% Part: normal-modal-logic
% Chapter: tableaux
% Section: countermodels

\documentclass[../../../include/open-logic-section]{subfiles}

\begin{document}

\olfileid{nml}{tab}{cou}

\olsection{\usetoken{P}{tableau} నుంచి ప్రతినమూనాలు}

\begin{explain}
  సంపూర్ణత సిద్ధాంతం నిరూపణ $\Entails !A$ అయితే
  $\Proves !A$ అని మాత్రమే చూపదు; $\Entails/ !A$ అయితే
  $!A$కు ప్రతినమూనా నిర్మించే పద్ధతినీ ఇస్తుంది.
  $\Log{K}$ సందర్భంలో ఈ పద్ధతి ఒక \emph{నిర్ణయ విధానం}.
  $\Entails/ !A$ అనుకుందాం. అప్పుడు
  \olref[cpl]{prop:complete-tableau} నిరూపణ సంపూర్ణ
  \tetoken{టాబ్లో}{!!{tableau}} నిర్మించే పద్ధతిని ఇస్తుంది.
  ఆ పద్ధతి ఎప్పుడూ ముగుస్తుంది. $\Log{K}$ ప్రతిజ్ఞావాక్య
  నియమాలు తక్కువ సంక్లిష్టత గల పూర్వసూచికలతో కూడిన
  \tetoken{సూత్రాలనే}{!!{formula}s} చేర్చుతాయి; అంటే
  శాఖపై ప్రతి $\sFmla{S}{!A}[\sigma]$కు ప్రతిజ్ఞావాక్య
  నియమాన్ని ఒక్కసారే వర్తింపజేయాలి. కొత్త పూర్వసూచికలు
  కింది నియమాల వల్ల మాత్రమే వస్తాయి:
  \iftag{prvBox}{$\TRule{\False}{\Box}$}{}\iftag{notprvBox,notprvDiamond}{}{
    and }\iftag{prvDiamond}{$\TRule{\True}{\Diamond}$}{}
  \iftag{notprvBox,notprvDiamond}{నియమం}{నియమాలు}.
  వాటినీ ఒక్కసారే వర్తింపజేయాలి; ప్రతి వర్తింపు ఒక్క కొత్త
  పూర్వసూచికనే ఉత్పత్తి చేస్తుంది.
  \iftag{prvBox}{$\TRule{\True}{\Box}$}{}\iftag{notprvBox,notprvDiamond}{}{
    and }\iftag{prvDiamond}{$\TRule{\False}{\Diamond}$}{}
  \iftag{notprvBox,notprvDiamond}{నియమాన్ని}{నియమాలను}
  బహుశా పలుమార్లు వర్తింపజేయాలి; కానీ ప్రతి పూర్వసూచికకు
  ఒక్కసారే. కొత్త పూర్వసూచికలు పరిమిత సంఖ్యలోనే ఉత్పత్తి
  అవుతాయి. అందువల్ల పరిమిత దశల్లో నిర్మాణం సంవృత శాఖను
  లేదా సంపూర్ణ శాఖను ఇస్తుంది.
  \sourcecorrection{OLTENMLTABCM-001}{మొదటి ప్రతినమూనా
    వాక్యంలో మూలం $\Entails/ A$ అని ఆశ్చర్యసూచికను
    వదిలింది; ఈ పేరాలోని మిగిలిన వాక్యాల ప్రకారం
    $\Entails/ !A$గా పునరుద్ధరించాం.}

  వివృత సంపూర్ణ శాఖ గల టాబ్లో
  నిర్మించాక, \olref[cpl]{thm:tableau-completeness} నిరూపణ
  అసలు పూర్వసూచికలతో కూడిన \tetoken{సూత్రాల}{!!{formula}s}
  సమితిని సంతృప్తిపరిచే స్పష్టమైన నమూనాను ఇస్తుంది.
  కాబట్టి $\Gamma \Entails !A$ అయితే సంవృత
  \tetoken{టాబ్లో}{!!{tableau}} ఉండి $\Gamma \Proves !A$
  అవుతుందని మాత్రమే కాదు; సంవృత
  \tetoken{టాబ్లో}{!!{tableau}} కోసం సరైన పద్ధతిలో
  వెతికి చివరికి ``సంపూర్ణ''
  \tetoken{టాబ్లో}{!!{tableau}} వస్తే,
  $\Gamma \Entails/ !A$ అని తెలుసుకోవడంతో పాటు
  ప్రతినమూనానూ నిర్మించగలం.
\end{explain}

\iftag{prvBox}{
\begin{ex}
  $\Proves/ \Box(p \lor q) \lif (\Box p \lor \Box
  q)$ అని మనకు తెలుసు. టాబ్లో
  నిర్మాణం ఇలా మొదలవుతుంది:
  \begin{oltableau}
    [\pFmla{\False}{\Box(p \lor q) \lif (\Box p \lor \Box q)}{1},
      just = \TAss, checked
      [\pFmla{\True}{\Box(p \lor q)}{1},
        just = {\TRule{\False}{\lif}[1]},
        [\pFmla{\False}{\Box p \lor \Box q}{1},
          just = {\TRule{\False}{\lif}[1]}, checked
          [\pFmla{\False}{\Box p}{1},
            just = {\TRule{\False}{\lor}[3]}, checked
            [\pFmla{\False}{\Box q}{1},
              just = {\TRule{\False}{\lor}[3]}, checked
              [\pFmla{\False}{p}{1.1},
                just = {\TRule{\False}{\Box}[4]}, checked
                [\pFmla{\False}{q}{1.2},
                  just = {\TRule{\False}{\Box}[5]}, checked
                ]
              ]
            ]
          ]
        ]
      ]
    ]
  \end{oltableau}
  ఈ \tetoken{టాబ్లో}{!!{tableau}} ఇంకా పూర్తి కాలేదు.
  తరువాత టిక్ గుర్తు లేని ఏకైక పంక్తిని, అంటే
  $2$వ పంక్తిలోని పూర్వసూచిక గల
  \tetoken{సూత్రాన్ని}{!!{formula}}
  $\sFmla{\True}{\Box(p \lor q)}[1]$ పరిశీలిద్దాం.
  శాఖపై ఉపయోగించిన ప్రతి పూర్వసూచికకు, అంటే
  $1.1$, $1.2$ రెండింటికీ $\TRule{\True}{\Box}$
  నియమాన్ని వర్తింపజేస్తాం:
  \begin{oltableau}
    [\pFmla{\False}{\Box(p \lor q) \lif (\Box p \lor \Box q)}{1},
      just = \TAss, checked
      [\pFmla{\True}{\Box(p \lor q)}{1},
        just = {\TRule{\False}{\lif}[1]},
        [\pFmla{\False}{\Box p \lor \Box q}{1},
          just = {\TRule{\False}{\lif}[1]}, checked
          [\pFmla{\False}{\Box p}{1},
            just = {\TRule{\False}{\lor}[3]}, checked
            [\pFmla{\False}{\Box q}{1},
              just = {\TRule{\False}{\lor}[3]}, checked
              [\pFmla{\False}{p}{1.1},
                just = {\TRule{\False}{\Box}[4]}, checked
                [\pFmla{\False}{q}{1.2},
                  just = {\TRule{\False}{\Box}[5]}, checked
                  [\pFmla{\True}{p \lor q}{1.1},
                    just = {\TRule{\True}{\Box}[2]}
                    [\pFmla{\True}{p \lor q}{1.2},
                      just = {\TRule{\True}{\Box}[2]}
                    ]
                  ]
                ]
              ]
            ]
          ]
        ]
      ]
    ]
  \end{oltableau}
  ఇప్పుడు 2, 8, 9 పంక్తులకు టిక్ గుర్తులు లేవు.
  అయితే కొత్త పూర్వసూచిక చేరలేదు; కాబట్టి 8, 9
  పంక్తులకు $\TRule{\True}{\lor}$ నియమాన్ని,
  శాఖ సంవృతమయ్యే వరకు, వచ్చే ప్రతి శాఖలో వర్తింపజేస్తాం:
  \begin{oltableau}
    [\pFmla{\False}{\Box(p \lor q) \lif (\Box p \lor \Box q)}{1},
      just = \TAss, checked
      [\pFmla{\True}{\Box(p \lor q)}{1},
        just = {\TRule{\False}{\lif}[1]}, checked
        [\pFmla{\False}{\Box p \lor \Box q}{1},
          just = {\TRule{\False}{\lif}[1]}, checked
          [\pFmla{\False}{\Box p}{1},
            just = {\TRule{\False}{\lor}[3]}, checked
            [\pFmla{\False}{\Box q}{1},
              just = {\TRule{\False}{\lor}[3]}, checked
              [\pFmla{\False}{p}{1.1},
                just = {\TRule{\False}{\Box}[4]}, checked
                [\pFmla{\False}{q}{1.2},
                  just = {\TRule{\False}{\Box}[5]}, checked
                  [\pFmla{\True}{p \lor q}{1.1},
                    just = {\TRule{\True}{\Box}[2]}, checked
                    [\pFmla{\True}{p \lor q}{1.2},
                      just = {\TRule{\True}{\Box}[2]}, checked
                      [\pFmla{\True}{p}{1.1},
                        just = {\TRule{\True}{\lor}[8]}, checked, close
                      ]
                      [\pFmla{\True}{q}{1.1},
                        just = {\TRule{\True}{\lor}[8]}, checked
                        [\pFmla{\True}{p}{1.2},
                          just = {\TRule{\True}{\lor}[9]}, checked]
                        [\pFmla{\True}{q}{1.2},
                          just = {\TRule{\True}{\lor}[9]}, checked, close]
                      ]
                    ]
                  ]
                ]
              ]
            ]
          ]
        ]
      ]
    ]
  \end{oltableau}
  ఒక్క వివృత శాఖ మిగులుతుంది; అది సంపూర్ణం. దాని నుంచి
  లోకాలు $W = \{1, 1.1, 1.2\}$ (వివృత శాఖపై కనిపించే
  పూర్వసూచికలే), ప్రాప్యత సంబంధం
  $R = \{\tuple{1, 1.1}, \tuple{1, 1.2}\}$,
  అర్థనిర్దేశం $V(p) = \{1.2\}$ (12వ పంక్తిలో
  $\sFmla{\True}{p}[1.2]$ ఉంది కాబట్టి),
  $V(q) = \{1.1\}$ (11వ పంక్తిలో
  $\sFmla{\True}{q}[1.1]$ ఉంది కాబట్టి) గల
  నమూనాను నిర్వచిద్దాం.
  \sourcecorrection{OLTENMLTABCM-002}{మొదటి పూర్తి
    టాబ్లోలో $\sFmla{\True}{p}[1.2]$ పంక్తి 12,
    $\sFmla{\True}{q}[1.1]$ పంక్తి 11;
    మూల గద్యంలో వరుసగా 11, 10 అని తప్పుగా ఉంది.}
  నమూనాను \olref{fig:counter-Box}లో చూపాం;
  అది $\Box(p \lor q) \lif (\Box p \lor \Box q)$కు
  ప్రతినమూనా అని తనిఖీ చేయవచ్చు.
  \begin{figure}
  \begin{center}
    \begin{tikzpicture}[modal]
      \node[world] (w1) [label={[align=right]right:\mFalse{p}\\ \mFalse{q}}]{$1$};
      \node[world] (w2) [label={[align=right]right:\mFalse{p}\\ \mTrue{q}},
        above left=of w1]{$1.1$};
      \node[world] (w3) [label={[align=right]right:\mTrue{p}\\ \mFalse{q}},
        above right=of w1] {$1.2$};
      \draw[->] (w1) to (w2);
      \draw[->] (w1) to (w3);
    \end{tikzpicture}
  \end{center}
  \caption{$\Box(p \lor q) \lif (\Box p \lor \Box q)$కు
    ప్రతినమూనా.}
  \ollabel{fig:counter-Box}
\end{figure}
\end{ex}
}
{
\begin{ex}
  $\Proves/ (\Diamond p \land \Diamond q) \lif \Diamond(p
  \land q)$ అని మనకు తెలుసు. టాబ్లో
  నిర్మాణం ఇలా మొదలవుతుంది:
  \begin{oltableau}
    [\pFmla{\False}{(\Diamond p \land \Diamond q) \lif \Diamond(p \land q)}{1},
      just = \TAss, checked
      [\pFmla{\True}{\Diamond p \land \Diamond q}{1},
        just = {\TRule{\False}{\lif}[1]}, checked
        [\pFmla{\False}{\Diamond(p \land q)}{1},
          just = {\TRule{\False}{\lif}[1]}
          [\pFmla{\True}{\Diamond p}{1},
            just = {\TRule{\True}{\land}[2]}, checked
            [\pFmla{\True}{\Diamond q}{1},
              just = {\TRule{\True}{\land}[2]}, checked
              [\pFmla{\True}{p}{1.1},
                just = {\TRule{\True}{\Diamond}[4]}, checked
                [\pFmla{\True}{q}{1.2},
                  just = {\TRule{\True}{\Diamond}[5]}, checked
                ]
              ]
            ]
          ]
        ]
      ]
    ]
  \end{oltableau}
  ఈ \tetoken{టాబ్లో}{!!{tableau}} ఇంకా పూర్తి కాలేదు.
  తరువాత టిక్ గుర్తు లేని ఏకైక పంక్తిని, అంటే
  $3$వ పంక్తిలోని పూర్వసూచిక గల
  \tetoken{సూత్రాన్ని}{!!{formula}}
  $\sFmla{\False}{\Diamond(p \land q)}[1]$
  పరిశీలిద్దాం. శాఖపై ఉపయోగించిన ప్రతి పూర్వసూచికకు,
  అంటే $1.1$, $1.2$ రెండింటికీ
  $\TRule{\False}{\Diamond}$ నియమాన్ని వర్తింపజేస్తాం:
  \sourcecorrection{OLTENMLTABCM-003}{3వ పంక్తి మూల
    చెట్టులో $\False\Diamond$ అయినా మూల వివరణ
    $\True\Diamond$ అని, దాని నియమాన్నే పేర్కొంది;
    $\False\Diamond$గా సరిచేశాం.}
  \begin{oltableau}
    [\pFmla{\False}{(\Diamond p \land \Diamond q) \lif \Diamond(p \land q)}{1},
      just = \TAss, checked
      [\pFmla{\True}{\Diamond p \land \Diamond q}{1},
        just = {\TRule{\False}{\lif}[1]}, checked
        [\pFmla{\False}{\Diamond (p \land q)}{1},
          just = {\TRule{\False}{\lif}[1]}
          [\pFmla{\True}{\Diamond p}{1},
            just = {\TRule{\True}{\land}[2]}, checked
            [\pFmla{\True}{\Diamond q}{1},
              just = {\TRule{\True}{\land}[2]}, checked
              [\pFmla{\True}{p}{1.1},
                just = {\TRule{\True}{\Diamond}[4]}, checked
                [\pFmla{\True}{q}{1.2},
                  just = {\TRule{\True}{\Diamond}[5]}, checked
                  [\pFmla{\False}{p \land q}{1.1},
                    just = {\TRule{\False}{\Diamond}[3]}
                    [\pFmla{\False}{p \land q}{1.2},
                      just = {\TRule{\False}{\Diamond}[3]}
                    ]
                  ]
                ]
              ]
            ]
          ]
        ]
      ]
    ]
  \end{oltableau}
  \sourcecorrection{OLTENMLTABCM-004}{రెండో వజ్ర
    టాబ్లోలో మొదటి పంక్తి మూలంలోని ముందు/వెనుక భాగాలను
    తారుమారు చేసింది; మిగిలిన రెండు చెట్లు, వాటి
    $\False\lif$ విస్తరణలకు సరిపడే అసలు సూత్రాన్ని
    పునరుద్ధరించాం.}
  ఇప్పుడు 3, 8, 9 పంక్తులకు టిక్ గుర్తులు లేవు.
  కొత్త పూర్వసూచిక చేరలేదు కాబట్టి 8, 9 పంక్తులకు
  $\TRule{\False}{\land}$ నియమాన్ని, శాఖ సంవృతమయ్యే
  వరకు, వచ్చే ప్రతి శాఖలో వర్తింపజేస్తాం:
  \begin{oltableau}
    [\pFmla{\False}{(\Diamond p \land \Diamond q) \lif \Diamond(p \land q)}{1},
      just = \TAss, checked
      [\pFmla{\True}{\Diamond p \land \Diamond q}{1},
        just = {\TRule{\False}{\lif}[1]}, checked
        [\pFmla{\False}{\Diamond(p \land q)}{1},
          just = {\TRule{\False}{\lif}[1]}
          [\pFmla{\True}{\Diamond p}{1},
            just = {\TRule{\True}{\land}[2]}, checked
            [\pFmla{\True}{\Diamond q}{1},
              just = {\TRule{\True}{\land}[2]}, checked
              [\pFmla{\True}{p}{1.1},
                just = {\TRule{\True}{\Diamond}[4]}, checked
                [\pFmla{\True}{q}{1.2},
                  just = {\TRule{\True}{\Diamond}[5]}, checked
                  [\pFmla{\False}{p \land q}{1.1},
                    just = {\TRule{\False}{\Diamond}[3]}, checked
                    [\pFmla{\False}{p \land q}{1.2},
                      just = {\TRule{\False}{\Diamond}[3]}, checked
                      [\pFmla{\False}{p}{1.1},
                        just = {\TRule{\False}{\land}[8]}, checked, close
                      ]
                      [\pFmla{\False}{q}{1.1},
                        just = {\TRule{\False}{\land}[8]}, checked
                        [\pFmla{\False}{p}{1.2},
                          just = {\TRule{\False}{\land}[9]}, checked]
                        [\pFmla{\False}{q}{1.2},
                          just = {\TRule{\False}{\land}[9]}, checked, close]
                      ]
                    ]
                  ]
                ]
              ]
            ]
          ]
        ]
      ]
    ]
  \end{oltableau}
  ఒక్క వివృత శాఖ మిగులుతుంది; అది సంపూర్ణం. దాని నుంచి
  లోకాలు $W = \{1, 1.1, 1.2\}$ (వివృత శాఖపై కనిపించే
  పూర్వసూచికలే), ప్రాప్యత సంబంధం
  $R = \{\tuple{1, 1.1}, \tuple{1, 1.2}\}$,
  అర్థనిర్దేశం $V(p) = \{1.1\}$ (6వ పంక్తిలో
  $\sFmla{\True}{p}[1.1]$ ఉంది కాబట్టి),
  $V(q) = \{1.2\}$ (7వ పంక్తిలో
  $\sFmla{\True}{q}[1.2]$ ఉంది కాబట్టి) గల
  నమూనాను నిర్వచిద్దాం.
  \sourcecorrection{OLTENMLTABCM-005}{7వ పంక్తిలోని
    $\sFmla{\True}{q}$కు మూల వివరణలో $1.1$ అని
    ఉన్నా, చెట్టులోనూ నమూనా చిత్రంలోనూ అది $1.2$;
    ఆ పూర్వసూచికను సరిచేశాం.}
  నమూనాను \olref{fig:counter-Diamond}లో చూపాం;
  అది $(\Diamond p \land \Diamond q) \lif \Diamond (p \land q)$కు
  ప్రతినమూనా అని తనిఖీ చేయవచ్చు.
  \begin{figure}
  \begin{center}
    \begin{tikzpicture}[modal]
      \node[world] (w1) [label={[align=right]right:\mFalse{p}\\ \mFalse{q}}]{$1$};
      \node[world] (w2) [label={[align=right]right:\mTrue{p}\\ \mFalse{q}},
        above left=of w1]{$1.1$};
      \node[world] (w3) [label={[align=right]right:\mFalse{p}\\ \mTrue{q}},
        above right=of w1] {$1.2$};
      \draw[->] (w1) to (w2);
      \draw[->] (w1) to (w3);
    \end{tikzpicture}
  \end{center}
  \caption{$(\Diamond p \land \Diamond q) \lif
    \Diamond (p \land q)$కు ప్రతినమూనా.}
  \ollabel{fig:counter-Diamond}
\end{figure}
\end{ex}
}

\end{document}

